% Zdroj: PDF priklady-podzim-2024.pdf, fyzická str. 2--3. Značka: společné okruhy (matematika). Zařazeno: M1.
\section{Integrály a primitivní funkce}
\szzotazka{podzim 2024}{3}{Integrály a primitivní funkce}

\noindent\textbf{Zadání.}\par
\begin{enumerate}
  \item Definujte, co je \uv{primitivní funkce} k dané funkci $f$ na daném otevřeném intervalu $I$.
  \item Zformulujte pravidlo \uv{per partes} (též známé jako integrování po částech) pro výpočet primitivní funkce. Můžete se pro jednoduchost omezit na situaci, kdy všechny uvažované funkce jsou spojité.
  \item Nechť $P \subseteq \mathbb{R}^2$ je oblast roviny ohraničená zdola osou $x$ a shora grafem funkce $f(x) = x \cos(x)$ pro $x \in [0, \pi/2]$. Formálněji řečeno,
    \[
      P = \left\{ (x, y) \in \mathbb{R}^2 ;\ 0 \leq x \leq \frac{\pi}{2} \,\wedge\, 0 \leq y \leq x \cos(x) \right\}.
    \]
    Spočítejte plošný obsah oblasti $P$.
\end{enumerate}

\medskip
\noindent\textbf{Náčrt řešení.}\par
\begin{enumerate}
  \item Funkce $F$ je primitivní funkce k funkci $f$ na intervalu $I$, pokud má $F$ v každém bodě $x \in I$ derivaci rovnou $f(x)$.
  \item Pravidlo pro integrování per partes říká, že pokud na nějakém intervalu $I$ jsou $f$ a $g$ spojité funkce, $F$ je primitivní funkce k $f$ a $G$ je primitivní funkce ke $g$, pak na tomto intervalu platí
    \[
      \int f(x) G(x)\,\mathrm{d}x = F(x) G(x) - \int F(x) g(x)\,\mathrm{d}x.
    \]
  \item Hledaný plošný obsah lze vyjádřit pomocí integrálu jako $\int_0^{\pi/2} x \cos(x)\,\mathrm{d}x$. Hodnotu integrálu určíme použitím metody per partes pro výpočet určitého integrálu:
    \begin{align*}
      \int_0^{\pi/2} x \cos(x)\,\mathrm{d}x
        &= \bigl[x \sin(x)\bigr]_0^{\pi/2} - \int_0^{\pi/2} \sin(x)\,\mathrm{d}x \\
        &= \left(\frac{\pi}{2} - 0\right) - \bigl[-\cos(x)\bigr]_0^{\pi/2} \\
        &= \frac{\pi}{2} - 1.
    \end{align*}
\end{enumerate}
